\documentclass[tikz,border=10pt]{standalone}
\usepackage{ctex}
\usepackage{amsmath}
\usetikzlibrary{calc, arrows.meta, shapes, positioning}

\begin{document}
% 综合法 vs 向量法 对比图（以求二面角为例）
\begin{tikzpicture}[
  font=\small,
  step/.style={draw, rounded corners=3pt, fill=#1!10, draw=#1!60,
               minimum width=3.8cm, minimum height=0.7cm, align=center, inner sep=4pt},
  arrow/.style={-Stealth, thick, #1},
]

% ─── 标题 ───────────────────────────────────────────
\node[font=\normalsize\bfseries] at (4.2, 8.0) {立体几何：综合法 vs 向量法（求二面角）};

% ─── 左列：综合法 ────────────────────────────────────
\node[font=\bfseries, blue!70] at (1.5, 7.2) {综合法};
\node[step=blue] (s1) at (1.5, 6.5) {作垂线 / 辅助平面};
\node[step=blue] (s2) at (1.5, 5.65) {找二面角的棱和半平面};
\node[step=blue] (s3) at (1.5, 4.8)  {在半平面内作垂线段};
\node[step=blue] (s4) at (1.5, 3.95) {用三角形计算角度};
\node[step=blue] (s5) at (1.5, 3.1)  {需要辅助线（技巧强）};

\draw[arrow=blue] (s1) -- (s2);
\draw[arrow=blue] (s2) -- (s3);
\draw[arrow=blue] (s3) -- (s4);
\draw[arrow=blue] (s4) -- (s5);

% 步骤数标注
\node[blue!60, font=\scriptsize] at (1.5, 2.55) {约 5 步，依赖作图直觉};

% ─── 右列：向量法 ────────────────────────────────────
\node[font=\bfseries, orange!80!black] at (6.9, 7.2) {向量法};
\node[step=orange] (v1) at (6.9, 6.5) {建立空间直角坐标系};
\node[step=orange] (v2) at (6.9, 5.65) {写出各顶点坐标};
\node[step=orange] (v3) at (6.9, 4.8) {求两平面法向量 $\vec{n}_1, \vec{n}_2$};
\node[step=orange] (v4) at (6.9, 3.95) {$\cos\theta = \dfrac{|\vec{n}_1\cdot\vec{n}_2|}{|\vec{n}_1||\vec{n}_2|}$};
\node[step=orange] (v5) at (6.9, 3.1) {无须辅助线，纯代数};

\draw[arrow=orange!80!black] (v1) -- (v2);
\draw[arrow=orange!80!black] (v2) -- (v3);
\draw[arrow=orange!80!black] (v3) -- (v4);
\draw[arrow=orange!80!black] (v4) -- (v5);

\node[orange!60!black, font=\scriptsize] at (6.9, 2.55) {约 4 步，格式固定，不易出错};

% ─── 中间对比线 ─────────────────────────────────────
\draw[gray!50, dashed] (4.2, 2.3) -- (4.2, 7.4);
\node[gray, font=\scriptsize, rotate=90] at (4.2, 4.85) {适用场景对比};

% ─── 底部选择建议 ───────────────────────────────────
\node[draw=gray!60, fill=gray!8, rounded corners=4pt,
      font=\scriptsize, align=center, inner sep=6pt,
      text width=7.5cm] at (4.2, 1.55)
  {选择建议：\\
   规则图形（正方体/棱柱）且辅助线明显 → 综合法更简洁\\
   含参/折叠/不规则图形 → 向量法更稳健};

\end{tikzpicture}
\end{document}
